% Generated draft: requires mathematical and source review.
\documentclass[11pt]{article}
\usepackage[margin=25mm]{geometry}
\usepackage{fontspec}
\setmainfont{Latin Modern Roman}
\usepackage{amsmath,amssymb,mathtools,unicode-math}
\setmathfont{Latin Modern Math}
\usepackage{enumitem,needspace,xurl,hyperref}
\setlength{\parindent}{0pt}
\setlength{\parskip}{5pt}
\setlength{\emergencystretch}{3em}
\newcommand{\N}{\mathbb{N}}
\newcommand{\Z}{\mathbb{Z}}
\newcommand{\Q}{\mathbb{Q}}
\newcommand{\R}{\mathbb{R}}
\newcommand{\C}{\mathbb{C}}
\newcommand{\F}{\mathbb{F}}
\newcommand{\A}{\mathbb{A}}
\newcommand{\PP}{\mathbb{P}}
\newcommand{\E}{\mathbb{E}}
\newcommand{\eps}{\varepsilon}
\newcommand{\dd}{\,\mathrm{d}}
\begin{document}
\section*{UnsolvedMath ID: 6700020}
\textbf{Source problem code:} catalogue entry
\subsection{Definitions and mathematical statement}
The source proposes relaxing an earlier ``isometry at infinity'' hypothesis to suitable asymptotic flatness together with negativity of a suitable energy at infinity. The earlier statement, its manifolds, curvature hypotheses, and desired conclusion are absent from this record.

Being isometric at infinity ordinarily means an exact isometry outside specified compact subsets. Asymptotic flatness instead involves coordinates on an end in which a metric approaches a Euclidean metric, with prescribed decay of coefficients and derivatives. The source supplies neither a reference end nor decay, dimension, regularity or completeness conditions. It also does not define the proposed energy functional $E_\infty$ or its sign convention.

The literal proposed sign is negativity, which would mean $E_\infty<0$ under a fixed convention. This is different from $E_\infty=0$ or $E_\infty\leq0$. The supplied summary invokes positive-mass and zero-mass rigidity results: in their appropriate complete asymptotically flat, nonnegative-scalar-curvature, dimension/spin and decay settings, ADM mass is nonnegative and its zero case has a flatness conclusion. Those benchmark statements do not by themselves identify the source's energy, explain its negative sign, or recover the missing antecedent theorem. ADM and Bartnik notions must likewise not be conflated.

Thus the mathematical target is an unspecified asymptotic strengthening of an earlier geometric assertion. No precise replacement hypothesis or rigidity conclusion can be stated until that context and energy definition are supplied.
\subsection{Short English statement}
Can exact agreement at infinity in an earlier scalar-curvature statement be weakened to asymptotic flatness plus a negative energy condition? The original theorem, asymptotics, and energy are missing; zero-mass rigidity is a distinct benchmark.
\subsection{Sources}
\begin{itemize}
\item Local UnsolvedMath registry, permanent ID 6700020, exact source code catalogue entry, scalar-curvature item [?20]; proposed by Misha Gromov, 2017. Registry status partially solved, classification PARTIAL-PROGRESS.
\item Misha Gromov, \emph{101 Questions, Problems and Conjectures around Scalar Curvature} (2017), printed tag [?20], occurrence 20, PDF page 24, \url{https://www.ihes.fr/~gromov/wp-content/uploads/2018/08/101-problemsOct1-2017.pdf}. Supplied normalized pdf-text extraction accessed July 29, 2026; public author-hosted incomplete PDF. No PDF was accessed here.
\item The supplied summary invokes positive-mass results of Schoen--Yau and Witten, ADM/Bartnik energies, and ongoing generalizations to other ends and higher codimension, without exact references or hypotheses. Its identification of this fragment with positive-mass rigidity is contextual interpretation, not a verified equivalence of statements.
\item Provenance flags current status and statement-reuse terms NEEDS\_REVIEW; the source marks the passage as a numbered open item.
\end{itemize}
\end{document}
